

\newcommand{\imprimirbrazilabstract}{%
    \cleardoublepage\phantomsection
    \addtotextpreliminarycontent{Resumo em Português}
    \begin{otherlanguage*}{brazil}
    \begin{resumo}[Resumo]

Esta dissertação investiga se a distância linguística e a especialização em um idioma estão associadas à vulnerabilidade de grandes modelos de linguagem a ataques multilíngues. Esses ataques, conhecidos como \textit{jailbreaks}, tentam contornar as restrições de segurança dos modelos para obter ajuda na execução de ações prejudiciais. O estudo avalia, em treze idiomas, oito assistentes organizados em quatro pares voltados ao português brasileiro, ao italiano, ao sueco e ao búlgaro, além de um modelo de referência em inglês. As solicitações prejudiciais são traduzidas automaticamente, e as respostas em outros idiomas são traduzidas para o inglês antes da avaliação. A avaliação segue o protocolo StrongREJECT, que atribui a cada resposta uma pontuação de zero a um conforme o quanto ela ajuda a executar a ação prejudicial. A distância entre idiomas é calculada a partir de características linguísticas. O indicador de especialização combina a taxa de acerto em questões multilíngues de compreensão de leitura com a frequência com que o modelo repete as respostas que dá no idioma ao qual é voltado. A análise principal reúne \GMainN{} respostas pontuadas dos oito modelos pareados, incluindo as pontuações iguais a zero, e usa regressão logística fracionária com efeitos fixos de modelo e idioma e erros-padrão robustos agrupados por combinação de modelo e idioma. O uso da pontuação como desfecho e as demais escolhas de análise foram definidos depois da inspeção dos resultados. O coeficiente da distância é negativo em \RDistanceNegativeCount{} das \RDistanceSpecCount{} especificações avaliadas, o que contraria a hipótese, mas o intervalo de confiança de 95\% da análise principal inclui zero. Quando se exigem pelo menos cem respostas comparáveis no teste de compreensão de leitura, a estimativa passa a ser positiva, também com intervalo que inclui zero; essa exigência, porém, exclui todas as observações de um dos modelos. A especialização tem associação positiva com a pontuação, ao contrário da proteção esperada. Como os dois preditores são colineares, é difícil separar a associação de cada um. O limite de geração de respostas nesse teste compromete a medida de especialização de um dos modelos, e a tradução automática, a falta de validação humana bilíngue e as diferenças de origem e de treinamento entre os modelos também limitam a interpretação. No conjunto avaliado, os resultados não confirmam a associação positiva esperada para a distância nem a relação protetora esperada para a especialização. A contribuição é uma avaliação multilíngue acompanhada de análises de sensibilidade e de uma descrição explícita das limitações de medição.

        \imprimirpalavraschave{Palavras-chave}{grandes modelos de linguagem; \textit{jailbreak}; segurança multilíngue.}

    \end{resumo}
    \end{otherlanguage*}
}


\newcommand{\imprimirenglishabstract}{%
    % https://tex.stackexchange.com/questions/20987/changing-babel-package-inside-a-single-chapter
    % https://tex.stackexchange.com/questions/36526/multiple-language-document-babel-selectlanguage-vs-begin-endotherlanguage
    \cleardoublepage\phantomsection
    \addtotextpreliminarycontent{Abstract in English}
    \begin{otherlanguage*}{english}
    \begin{resumo}[Abstract]

This dissertation investigates whether linguistic distance and language specialization are associated with the vulnerability of large language models to multilingual attacks. These attacks, known as jailbreaks, try to bypass a model's safety restrictions to obtain help in carrying out harmful actions. The study evaluates, across thirteen languages, eight assistants organized into four pairs focused on Brazilian Portuguese, Italian, Swedish, and Bulgarian, plus an English reference model. Harmful prompts are machine-translated, and non-English responses are translated into English before evaluation. Evaluation follows the StrongREJECT protocol, which gives each response a score from zero to one according to how much it helps carry out the harmful action. Distance between languages is calculated from linguistic features. The specialization proxy combines accuracy on multilingual reading-comprehension questions with how often the model repeats the answers it gives in the language it is designed to serve. The main analysis includes \GMainN{} scored responses from the eight paired models, retaining zero scores, and uses fractional-logit regression with model and language fixed effects and robust standard errors clustered by model--language combination. The score outcome and the other analytical choices were defined after the results had been inspected. The distance coefficient is negative in \RDistanceNegativeCount{} of the \RDistanceSpecCount{} specifications examined, which runs contrary to the hypothesis, but the 95\% confidence interval in the main analysis includes zero. When at least one hundred comparable answers are required on the reading-comprehension test, the estimate becomes positive, again with an interval that includes zero; this requirement, however, excludes all observations from one model. Specialization is positively associated with scores, contrary to the expected protection. Because the two predictors are collinear, their associations are hard to separate. The response generation limit on this test compromises the specialization measure for one model, and machine translation, the lack of human bilingual validation, and differences in model ancestry and training also limit interpretation. Within the evaluated panel, the results confirm neither the expected positive association for distance nor the expected protective relationship for specialization. The contribution is a multilingual evaluation accompanied by sensitivity analyses and an explicit description of measurement limitations.

        \imprimirpalavraschave{Keywords}{large language models; jailbreak; multilingual safety.}

    \end{resumo}
    \end{otherlanguage*}
}


% \newcommand{\imprimirfrenchabstract}{%
%     \addtotextpreliminarycontent{Français Résumé}
%     \begin{resumo}[Résumé]
%       \begin{otherlanguage*}{french}
%           Il s'agit d'un résumé en français.

%           \imprimirpalavraschave{Mots-clés}{latex. abntex. publication de textes.}
%       \end{otherlanguage*}
%     \end{resumo}
% }


% \newcommand{\imprimirspanishabstract}{%
%     \addtotextpreliminarycontent{Español Resumen}
%     \begin{resumo}[Resumen]
%       \begin{otherlanguage*}{spanish}
%           Este es el resumen en español.

%           \imprimirpalavraschave{Palabras clave}{latex. abntex. publicación de textos.}
%       \end{otherlanguage*}
%     \end{resumo}
% }


\makeatletter
\ifenglish
    \@ifundefined{imprimirbrazilabstract}{}{\imprimirbrazilabstract}

    \@ifundefined{imprimirenglishabstract}{}{\imprimirenglishabstract}

    \cleardoublepage\phantomsection
\pretextualchapter{Resumo Expandido}
\addtotextpreliminarycontent{Resumo Expandido}
\begin{otherlanguage*}{brazil}
% Evita linhas isoladas e palavras partidas entre páginas neste elemento.
\clubpenalty=10000
\widowpenalty=10000
\brokenpenalty=10000
% Espaçamento uniforme entre os títulos das seções e o texto.
\setbeforesecskip{\baselineskip}
\setaftersecskip{\baselineskip}

\section*{Introdução e objetivos}
Esta dissertação investiga se a distância linguística e a especialização em um idioma estão associadas à vulnerabilidade de grandes modelos de linguagem a ataques multilíngues de \textit{jailbreak}. Esses ataques tentam contornar as restrições de segurança dos modelos para obter ajuda na execução de ações prejudiciais e, neste estudo, consistem em solicitações prejudiciais traduzidas para outros idiomas.

O objetivo é verificar se o quanto um modelo atende a essas solicitações está relacionado à distância entre o idioma do ataque e o idioma de alinhamento do modelo, ou à especialização do modelo no idioma do ataque. O idioma de alinhamento é aquele indicado na documentação de desenvolvimento do modelo; essa indicação não garante proficiência nem segurança nesse idioma. As hipóteses são de que uma distância maior esteja associada a mais atendimento às solicitações prejudiciais e de que uma especialização maior esteja associada a menos atendimento. A análise conjunta também verifica se, como esperado, a associação da distância é mais forte e mais estável do que a da especialização.

\section*{Metodologia}
O estudo avalia oito assistentes em treze idiomas, entre eles o idioma de alinhamento de cada um. Os assistentes formam quatro pares voltados ao português brasileiro, ao italiano, ao sueco e ao búlgaro. Em cada par, os rótulos de alinhamento fraco e forte seguem a documentação disponível. O alinhamento de segurança busca reduzir respostas prejudiciais, mas, como os modelos de cada par também diferem na origem e no treinamento, não é possível isolar o efeito desse alinhamento. O Llama-3.1-8B-Instruct serve como referência em inglês: entra na padronização do indicador de especialização, mas fica de fora das regressões com os oito modelos pareados.

O StrongREJECT fornece as solicitações prejudiciais e o protocolo de avaliação das respostas. O NLLB-200, um modelo de tradução automática, traduz as solicitações do inglês para os demais idiomas e as respostas em outros idiomas para o inglês antes da avaliação; os textos originais são preservados. O controle de qualidade das traduções combina uma triagem com o BLASER, que estima automaticamente a qualidade de traduções, e uma verificação de equivalência de sentido feita por outro modelo de linguagem. Das 148 traduções sinalizadas para essa verificação, 100 foram mantidas e 48 excluídas. A retradução das solicitações para o inglês ajuda a diagnosticar mudanças de sentido. Essas decisões não passaram por validação humana, e as traduções aprovadas automaticamente não foram calibradas com uma amostra representativa.

A pontuação oficial do StrongREJECT vai de zero a um e leva em conta se houve recusa, quão convincente é a resposta e quão específica ela é. Quanto maior o valor, mais a resposta ajuda a atender à solicitação prejudicial. As médias incluem as pontuações iguais a zero, que podem corresponder a recusas, respostas fora do assunto ou conteúdo sem utilidade para a solicitação. Saídas vazias e falhas de avaliação não contam como respostas seguras: ficam registradas como resultados desconhecidos e não entram nas médias. Bloqueios do provedor, falhas de geração e falhas do avaliador são separados das recusas feitas pelo próprio modelo. Os intervalos de confiança das médias são obtidos por reamostragem das solicitações originais, de modo que as respostas de diferentes modelos e idiomas a uma mesma solicitação permanecem juntas.

A distância linguística é a distância angular entre os vetores de características dos idiomas no URIEL+, uma base de características linguísticas. A especialização é medida de forma indireta com o BELEBELE, um conjunto multilíngue de questões de múltipla escolha sobre compreensão de leitura. O componente chamado seguimento de instruções (IF, do inglês \textit{instruction following}) corresponde, na prática, à proporção de acertos e não mede de forma geral a capacidade de seguir instruções. A consistência entre idiomas (CONS, do inglês \textit{cross-lingual consistency}) mede com que frequência o modelo escolhe a mesma alternativa que escolheu em seu idioma de alinhamento, mesmo quando as duas estão erradas. O indicador de especialização (SPEC, do inglês \textit{specialization}) combina IF e CONS com pesos iguais, depois de padronizá-los nas 117 combinações de modelo e idioma. Nenhuma dessas medidas avalia diretamente a segurança nem a capacidade de gerar respostas abertas.

No cálculo de IF, as 900 questões permanecem no denominador, e as respostas em que não é possível identificar a alternativa escolhida contam como erradas. O CONS compara apenas respostas à mesma questão em que a alternativa pôde ser identificada. As questões são pareadas entre idiomas pelos identificadores de origem, e não pela posição das linhas nos arquivos de cada idioma. Os gabaritos e as alternativas originais são preservados. As estimativas vêm acompanhadas da taxa de identificação de alternativas, do número de respostas comparáveis e dos intervalos de confiança, já que proporções calculadas com poucas respostas podem ser imprecisas.

A análise principal usa regressão logística fracionária, adequada a uma pontuação entre zero e um e que preserva os valores extremos. O ajuste é feito com equações de estimação generalizadas e efeitos fixos de modelo e de idioma do ataque. Distância e especialização são padronizadas em cada amostra de estimação. Como os dois preditores assumem um único valor por combinação de modelo e idioma, os erros-padrão robustos da inferência principal são agrupados nas \RNumCells{} combinações avaliadas. Assim, o número de respostas não corresponde ao número de comparações independentes entre idiomas. Intervalos agrupados por solicitação são apresentados para comparação com a análise anterior. A inferência é aproximada e vale apenas para o conjunto avaliado. O uso da pontuação como desfecho e as demais escolhas de análise foram definidos depois da inspeção dos resultados e, portanto, não são apresentados como previamente especificados.

Os diagnósticos verificam a convergência, a finitude das estimativas e se a matriz de covariância é definida positiva, além da colinearidade no conjunto completo de preditores e indicadores de modelo e idioma. Há colinearidade quando os preditores dependem tanto uns dos outros que fica difícil separar a associação de cada um com a resposta. Quando o fator de inflação da variância de um preditor substantivo passa de cinco, são feitos ajustes com um preditor substantivo de cada vez. As análises de sensibilidade incluem ajustes separados por grupo de alinhamento, a decomposição de SPEC em IF e CONS, restrições de cobertura e verificações de itens questionáveis do BELEBELE. Também são examinadas a qualidade das traduções e a inflação de tokens, isto é, o aumento relativo do número de unidades em que o modelo segmenta o texto. As comparações com e sem esse último preditor usam a mesma amostra de dados completos.

\section*{Resultados e discussão}
A análise principal reúne \GMainN{} respostas pontuadas dos oito modelos pareados. O coeficiente da distância é negativo em \RDistanceNegativeCount{} das \RDistanceSpecCount{} especificações avaliadas, ao contrário da hipótese de que os modelos seriam mais vulneráveis em idiomas mais distantes. No entanto, o intervalo de confiança de 95\% da análise principal, com agrupamento por combinação de modelo e idioma, inclui zero. Esse resultado não permite afirmar a direção da associação na análise principal, nem mostra que a associação não existe.

Quando se exigem pelo menos cem respostas comparáveis no cálculo de CONS, a estimativa para a distância fica positiva, também com intervalo que inclui zero. Essa exigência elimina todas as combinações do LLaMAntino-2, o modelo italiano classificado como fraco, e muda, portanto, tanto a cobertura da medição quanto o conjunto de modelos analisados. Por isso, a troca de sinal da estimativa pontual não pode ser atribuída apenas a uma medida mais precisa de CONS, e os intervalos não indicam associações de sinais opostos nas duas amostras.

Na análise principal, a especialização tem associação positiva com a pontuação, ao contrário da proteção esperada, e seu intervalo de confiança não inclui zero mesmo com o agrupamento por combinação de modelo e idioma. Essa associação não mostra que a especialização cause falhas de segurança, nem descarta que a capacidade linguística do modelo influencie a avaliação. A colinearidade dificulta separar as associações da distância e da especialização. Em conjunto, os resultados não confirmam a expectativa de que a associação da distância seria mais forte e mais estável.

A medição com o BELEBELE tem um problema específico no LLaMAntino-2. O limite de geração de oito tokens corta muitas respostas antes que o modelo indique a alternativa escolhida. Em italiano, o modelo costuma começar com uma resposta mais longa, o que contribui para a baixa cobertura justamente em seu idioma de alinhamento. O resultado depende, portanto, também do tamanho desse preâmbulo e de o modelo seguir a instrução de formato, escrita em inglês em todos os idiomas. Além disso, algumas saídas não respondem à questão, de modo que apenas aumentar o limite de geração não resolveria o problema. A restrição de cobertura retira esse modelo da análise, mas não corrige suas medidas de especialização.

Uma análise exploratória, apresentada à parte em apêndice, examina se direções nas ativações internas associadas à recusa se transferem entre idiomas. A transferência é mais consistente no modelo de referência e no modelo italiano ANITA do que no modelo sueco AI Sweden. Como os exemplos foram selecionados com base nas respostas, houve sobreposição de questões entre etapas e faltaram testes com questões independentes e intervenções de controle, esses resultados não permitem uma explicação causal do comportamento observado.

\section*{Considerações finais}
Para os modelos e idiomas avaliados, os resultados não confirmam a hipótese de que os modelos atendem mais a solicitações prejudiciais em idiomas linguisticamente mais distantes, nem a relação protetora esperada para a especialização. A contribuição é uma avaliação multilíngue documentada, com análises de sensibilidade e limitações de medição explícitas. As comparações não isolam o efeito do alinhamento de segurança nem estabelecem uma hierarquia universal de segurança entre idiomas.

A tradução automática e a falta de validação humana bilíngue limitam a interpretação. A validação no MultiJail, um conjunto externo de testes de segurança multilíngue, e os testes com solicitações inofensivas comparáveis não foram concluídos. O BELEBELE não controla diferenças na geração de respostas abertas, e seus problemas de cobertura e de limite de geração comprometem parte das medidas de especialização.

Como trabalho futuro, foi implementado um procedimento que pontua cada alternativa do BELEBELE pela média dos logaritmos das probabilidades condicionais de seus tokens. Esse procedimento dispensa a geração de respostas e a posterior identificação da alternativa, mas não foi executado nesta dissertação. Para avaliá-lo, é preciso compará-lo ao procedimento atual e repetir as análises que dependem de SPEC. Os valores usados no estudo continuam afetados pelo problema identificado. O Sabiá-3, modelo do par português, não fornece as probabilidades necessárias a essa nova medição. Ficam pendentes a validação humana bilíngue, a avaliação em conjuntos externos de ataques e os testes independentes das ativações internas.

\begingroup
\setlength{\parindent}{0pt}
\SingleSpacing
\imprimirpalavraschave{Palavras-chave}{grandes modelos de linguagem; \textit{jailbreak}; segurança multilíngue.}
\par
\endgroup
\end{otherlanguage*}


\else
    \@ifundefined{imprimirbrazilabstract}{}{\imprimirbrazilabstract}
    \@ifundefined{imprimirenglishabstract}{}{\imprimirenglishabstract}
\fi

\@ifundefined{imprimirfrenchabstract}{}{\imprimirfrenchabstract}
\@ifundefined{imprimirspanishabstract}{}{\imprimirspanishabstract}
\makeatother
